%% 第8巻 巻末（こたえ・記号カード・積分の地図・定数と単位・公式・チェック表・つぎの巻へ・さくいん・奥付）（第14回編集委員会）
%% 第6巻・第7巻の巻末と同じ形。記号カードの一覧・書くときのちゅうい・単位と数・公式・表計算と Python・専門書・つぎの巻へを、
%% 第8巻の本文（8-0〜8-14）に合わせて書いた。ルビなし。字の色は \bkTT。つぎの巻へのカードは第9巻の色（bkVol9・bkDeep9）。
%% 委員会の決定10：依頼者の指示（10月29日 13:34）を受けて、よく使う公式の節の最初に「積分の地図」（見開き2ページ）を置く。
%%   型は tools/kata8-int.md（組 V8-4）。場合の記号と例題の番号は、場合分けの表8-6・8-7・8-8・8-11 と各章の例題から写した。
%%   見開きにするため、地図の前で偶数ページ（左）にそろえる（\bkevenpage）。左＝見分け方・どの表を見る？（流れ図）・表8-6、右＝表8-7・表8-8・表8-11。
%% 記号は委員会の決定どおり：∮（閉じた道）・∯（閉じた面）、∬・∭、grad・div・rot と ∇、∇²、f_x を使わない、\vec n\,dS、電位 φ と球座標 φ（varphi）、k・α。
%% 数と式は本文と同じ（scratchpad/back8/check.py で確かめた）。
\makeatletter
%% ---------------- この巻末だけで使う部品（頭文字 kviiiz：kata8-int.md の型） ----------------
\colorlet{kviiizNc}{bkOrange!85!black}% 名前の枠の線（点線）
\colorlet{kviiizSc}{bkTheme}% ふつうの定積分の枠の線（実線）
\colorlet{kviiizSb}{bkTheme!8}% ふつうの定積分の地
\tikzset{
  kviiiz dan/.style={font=\sffamily\bfseries\footnotesize, rounded corners=1.2mm, inner xsep=1.6mm, inner ysep=0.5mm, line width=0.9pt},
  kviiiz dan1/.style={kviiiz dan, draw=kviiizNc, dash pattern=on 2.2pt off 1.4pt, fill=white, text=bkInk},
  kviiiz dan2/.style={kviiiz dan, draw=kviiizSc, fill=kviiizSb, text=bkInk},
  kviiiz dan3/.style={kviiiz dan, draw=bkTheme, fill=bkTheme, text=bkOn},
  kviiiz name/.style={draw=kviiizNc, line width=0.9pt, dash pattern=on 2.2pt off 1.4pt, rounded corners=1.2mm, inner xsep=1mm, inner ysep=0.9mm, fill=white},
  kviiiz seki/.style={draw=kviiizSc, line width=0.9pt, rounded corners=1.2mm, inner xsep=1mm, inner ysep=0.9mm, fill=kviiizSb},
  kviiiz tag/.style={font=\sffamily\bfseries\footnotesize, text=bkInk, inner sep=0.4mm},
  kviiiz q/.style={diamond, aspect=3.6, draw=bkYellow!75!black, line width=0.9pt, fill=bkYellowBg, font=\sffamily\footnotesize, align=center, inner sep=0.2mm, text=bkInk},
  kviiiz do/.style={rectangle, draw=\bkTT, line width=0.9pt, fill=white, rounded corners=0.6mm, font=\sffamily\footnotesize, align=center, inner xsep=1.5mm, inner ysep=1.1mm, text=bkInk},
  kviiiz go/.style={-{Stealth[length=2.2mm]}, line width=0.9pt, bkInk!75, shorten <=1.5mm, shorten >=1.5mm},
  kviiiz yn/.style={font=\sffamily\bfseries\footnotesize, text=bkInk, fill=white, inner sep=0.5mm},
}
% 段の札  \kviiizdan{1|2|3}
\newcommand\kviiizdan[1]{\tikz[baseline=(d.base)]\node[kviiiz dan#1] (d) {\vphantom{あ}\ifcase#1\or 名前\or 直す\or 楽をする\fi};}
% 場合の札（地図の表の中は小さめ）  \kviiizcase{C}
\newcommand\kviiizcase[1]{\tikz[baseline=(c.base)]\node[fill=bkTheme, text=bkOn, font=\sffamily\bfseries\footnotesize, rounded corners=0.7mm, inner xsep=0.8mm, inner ysep=0.2mm, minimum width=5mm] (c) {\vphantom{Ag}#1};}
% 表の名札（地図の帯の左）  \kviiiztab{8-6}
\newcommand\kviiiztab[1]{\tikz[baseline=(t.base)]\node[fill=bkTheme, text=bkOn, font=\sffamily\bfseries\footnotesize, rounded corners=0.8mm, inner xsep=1.2mm, inner ysep=0.6mm] (t) {\vphantom{Ag}表#1};}
% 式の中の「名前」（点線）と「ふつうの定積分」（実線）。地図では文字の大きさの式（\textstyle）
\newcommand\kviiizN[1]{\tikz[baseline=(n.base)]\node[kviiiz name] (n) {$#1$};}
\newcommand\kviiizS[1]{\tikz[baseline=(n.base)]\node[kviiiz seki] (n) {$#1$};}
\newcommand\kviiizND[1]{\tikz[baseline=(n.base)]\node[kviiiz name] (n) {$\displaystyle #1$};}
\newcommand\kviiizNt[1]{\tikz[baseline=(n.base)]{\node[kviiiz name] (n) {$\displaystyle #1$};\node[kviiiz tag, anchor=north] at ([yshift=-0.8mm]n.south) {名前};}}
\newcommand\kviiizSt[1]{\tikz[baseline=(n.base)]{\node[kviiiz seki] (n) {$\displaystyle #1$};\node[kviiiz tag, anchor=north] at ([yshift=-0.8mm]n.south) {ふつうの定積分};}}
% タイルの小さな絵（8-6 の \kviiigpic と同じ絵）  \kviiizpic{1〜7}
\newcommand\kviiizpic[1]{\begin{tikzpicture}[x=0.8mm, y=0.8mm, baseline=(current bounding box.center)]
  \ifcase#1\or
  \draw[bkInk, line width=0.6pt] (0,0) -- (14,0);
  \foreach \kviiizx/\kviiizh in {3/3.9,5/4.6,7/4.6,9/4.5,11/5.1} {\filldraw[fill=bkTheme!25, draw=bkTheme, line width=0.4pt] ({\kviiizx-1},0) rectangle ({\kviiizx+1},\kviiizh);}
  \draw[bkInk!80, line width=0.7pt] (1,3.2) .. controls (5,6) and (9,3) .. (13,5.6);
  \or
  \draw[bkGray, line width=0.4pt, step=2.5] (0,0) grid (12.5,7.5);
  \fill[bkTheme!45] (5,2.5) rectangle (7.5,5);
  \draw[bkInk, line width=0.7pt] (0,0) rectangle (12.5,7.5);
  \or
  \draw[bkInk, line width=0.6pt] (0,0) rectangle (9,6) (0,6) -- (3,8.5) -- (12,8.5) -- (9,6) (12,8.5) -- (12,2.5) -- (9,0);
  \filldraw[fill=bkTheme!45, draw=bkTheme, line width=0.4pt] (3.5,1.5) rectangle (5.5,3.5) (3.5,3.5) -- (4.3,4.2) -- (6.3,4.2) -- (5.5,3.5) (6.3,4.2) -- (6.3,2.2) -- (5.5,1.5);
  \or
  \draw[bkInk, line width=0.7pt] (0,1) .. controls (4,8) and (8,-1) .. (13,6);
  \draw[bkTheme, line width=1.8pt] (5.2,3.75) -- (7.4,2.9);
  \or
  \draw[bkInk, line width=0.7pt] (0,1) .. controls (4,8) and (8,-1) .. (13,6);
  \draw[bkTheme, line width=1.3pt, -{Stealth[length=1.8mm]}] (4.6,3.95) -- (7.8,2.85);
  \or
  \draw[bkInk, line width=0.6pt] (0,1) .. controls (4,4) and (8,4) .. (12,1) -- (13,6) .. controls (9,9) and (5,9) .. (1,6) -- cycle;
  \filldraw[fill=bkTheme!45, draw=bkTheme, line width=0.4pt] (5,3.6) -- (7.4,3.6) -- (7.8,5.8) -- (5.4,5.8) -- cycle;
  \or
  \draw[bkInk, line width=0.6pt] (0,1) .. controls (4,4) and (8,4) .. (12,1) -- (13,6) .. controls (9,9) and (5,9) .. (1,6) -- cycle;
  \filldraw[fill=bkTheme!45, draw=bkTheme, line width=0.4pt] (5,3.6) -- (7.4,3.6) -- (7.8,5.8) -- (5.4,5.8) -- cycle;
  \draw[bkRed, line width=1.1pt, -{Stealth[length=1.8mm]}] (6.4,4.7) -- (6.4,10.5);
  \fi
\end{tikzpicture}}
% 地図の表の字（9.5pt。行の間をつめる）
\def\kviiizfont{\fontsize{9.5}{12.5}\selectfont}
% 地図の表（列：記号／こんなとき→手順／例題。1行めは帯）
\NewTblrEnviron{kviiizmap}
\SetTblrInner[kviiizmap]{colspec={Q[c,m,wd=8.5mm] X[l,m] Q[c,m,wd=13mm]}, width=\linewidth, hlines={0.4pt,bkLine}, rowsep=1pt, colsep=1.5mm, stretch=0,
  cells={font=\kviiizfont}, row{1}={abovesep=1.6pt, belowsep=1.6pt}}
% 「こんなとき」（太字）と矢印  \kviiizif{こんなとき}
\newcommand\kviiizto{{\sffamily\bfseries\color{\bkTT}→}}
\newcommand\kviiizif[1]{{\sffamily\bfseries #1}\enspace\kviiizto\enspace}
% 地図の文（前書きと見分け方）
\def\kviiizbody{\fontsize{10}{15}\selectfont}
% 地図の表の帯（1行め。表のはば全体）  \kviiizband{中身}
\newcommand\kviiizband[1]{\parbox{136mm}{\raggedright\kviiizfont\vphantom{表}#1}}
% 地図の小見出し（太字の1行）
\newcommand\kviiizhead[1]{\par\addvspace{3pt}\noindent{\sffamily\bfseries\small\color{\bkTT}#1}\par\nobreak\vspace{1.5pt}}

%% ---------------- こたえ ----------------
\bkAnswerPages

%% ---------------- 記号カード ----------------
%% 前の巻の記号で、この巻でもよく使うもの（章の番号は、はじめて出た巻の章。この巻の章の \bkkigocard は自動でうしろに集まる）
%% 8-4 の grad のカードは記号を \fontsize{17}{19} で大きく見せている。一覧の記号の列（18 mm）では大きすぎるので、
%% 一覧の中だけ、17pt の指定をそのときの字の大きさ（\footnotesize）に読みかえる（章の本文のカードはそのまま。第6・7巻と同じ方法）
\let\bk@origfontsize\fontsize
\def\bk@fsseventeen{17}%
\protected\def\fontsize#1#2{\def\bk@fsarg{#1}\ifx\bk@fsarg\bk@fsseventeen\expandafter\bk@origfontsize\expandafter{\f@size}{\f@baselineskip}\else\bk@origfontsize{#1}{#2}\fi}%
%% 8-7 の線積分のカードは、大きなカードに合わせて \scalebox{0.62} で縮めている。一覧では小さすぎるので、一覧の中だけ 0.9 に読みかえる
\let\bk@origscalebox\scalebox
\def\bk@sbsmall{0.62}%
\protected\def\scalebox#1{\def\bk@sbarg{#1}\ifx\bk@sbarg\bk@sbsmall\def\bk@sbarg{0.9}\fi\expandafter\bk@origscalebox\expandafter{\bk@sbarg}}%
\bkKigoPages{%
\bkkigorow{$\Delta$}{デルタ}{変化の量。8-13 の $\Delta x$（マスのはば）・$\Delta t$（1歩の時間）。ラプラシアンの記号には使わない（8-12）}{2-3}%
\bkkigorow{$\dfrac{\dd y}{\dd x}$}{ディー・ワイ・ディー・エックス}{1つの変数の微分。変数が2つ以上の関数で1つだけ動かすときは $\partial$ で書く（8-2）}{5-3}%
\bkkigorow{$\approx$}{ニアリーイコール}{ほぼ等しい。全微分の近似 $\Delta f\approx\dd f$、誤差の見積もり（8-3）}{5-8}%
\bkkigorow{$\displaystyle\int_a^b$}{インテグラル}{ふつうの定積分。重積分・線積分・面積分は、これのくり返しに直して計算する（8-6）}{5-10}%
\bkkigorow{$\displaystyle\sum_{k=1}^{n}$}{シグマ}{全部たす。タイルの和、表計算の \texttt{SUM}（8-6・8-7）}{5-10}%
\bkkigorow{$e$}{イー}{ネイピア数。ガウス積分の $e^{-x^2}$（8-6）、熱方程式の答えの $e^{-\alpha a^2t}$（8-13）}{5-14}%
\bkkigorow{$\vec a$}{ベクトル・エー}{向きと大きさをもつ量。この巻では、場所ごとに矢印が決まる $\vec F(x,\ y)$（8-5）}{4-10}%
\bkkigorow{$|\vec a|$}{ぜったいち・エー}{矢印の長さ。$|\mathrm{grad}\,h|$ は坂の急さ（8-4）}{4-10}%
\bkkigorow{$\vec e$}{イー}{長さ1の矢印。$\vec e_x$・$\vec e_y$・$\vec e_z$、方向微分の $\vec e$（8-4）、$\vec e_r$（8-8）}{4-11}%
\bkkigorow{$\vec a\cdot\vec b$}{エー・ドット・ビー}{内積。方向微分（8-4）、仕事 $\vec F\cdot\dd\vec r$（8-7）、流束の $\vec F\cdot\vec n$（8-8）、$\nabla\cdot$（8-9）}{4-12}%
\bkkigorow{$\vec i,\ \vec j,\ \vec k$}{アイ・ジェイ・ケイ}{専門書の基本ベクトル。この本の $\vec e_x$・$\vec e_y$・$\vec e_z$ と同じ（8-4・8-5）}{4-13}%
\bkkigorow{$\vec a\times\vec b$}{エー・クロス・ビー}{外積。ななめの面の面積ベクトル $A\vec n$（8-8）、$\nabla\times$（8-10）}{4-14}%
\bkkigorow{$(r,\ \theta)$}{アール・シータ}{極座標。タイルの面積は $r\,\dd r\,\dd\theta$（8-3・8-6）}{4-15}%
\bkkigorow{\scalebox{0.86}{$(r,\ \theta,\ \varphi)$}}{アール・シータ・ファイ}{球座標。$\theta$ は真上から、$\varphi$ は床の上の角。$\dd V=r^2\sin\theta\,\dd r\,\dd\theta\,\dd\varphi$（8-6・8-8）}{4-15}%
\bkkigorow{$\phi$}{ファイ}{この巻では電位 $\phi$（8-4）。球座標の角 $\varphi$ とは別の文字}{5-15}%
\bkkigorow{$\omega$}{オメガ}{角速度〔rad/s〕。流れに置いた風車の角速度は $\frac12\,\mathrm{rot}\,\vec v$（8-10）}{4-8}%
\bkkigorow{$y_n$}{ワイ・エヌ}{$n$ 歩めの値。8-13 の $T_i^{\,n}$ は、場所の名札 $i$ もついた形}{6-2}%
\bkkigorow{$A^{\mathsf T}$}{エー・ティー}{転置。軸を回して見た矢印の成分は $R^{\mathsf T}\vec F$（8-5）}{7-1}%
\bkkigorow{$\det A$}{デターミナント・エー}{行列式。ヤコビアン $\det J$ は小さな面積の倍率（8-3・8-6）}{7-4}%
\bkkigorow{$\lambda$}{ラムダ}{固有値。8-13 の1歩の行列の固有値（こわれない条件）。熱伝導率には使わず $k$ と書く}{7-8}%
}
\let\fontsize\bk@origfontsize
\let\scalebox\bk@origscalebox

\def\bkbkhead#1{\par\addvspace{4mm}\noindent{\sffamily\bfseries\large\color{\bkTT}#1}\par\nobreak\vspace{2mm}}
%% 記号カードのつづき：手で書くとき・読むときのちゅうい（各章の「まちがえやすいところ」から）
\Needspace{8\baselineskip}
\bkbkhead{記号を書くときのちゅうい}
{\small 手で書くと、ほかの記号とまぎれやすいものと、読みまちがえやすい書き方です。積分の記号は、まず「名前」として読みます。}\par\vspace{2mm}
{\small
\begin{longtblr}{colspec={Q[c,m,wd=27mm] Q[c,m,wd=24mm] Q[l,m,wd=75mm]}, rowhead=1, row{1}={bg=bkVol!14, font=\sffamily\bfseries\footnotesize}, hlines={0.4pt,bkLine}, rowsep=0.8pt}
記号 & まぎれやすい & 書くとき・読むときのコツ \\
$h(20,\ 30)$ & $h(30,\ 20)$ & かっこの1つめが $x$（東）、2つめが $y$（北）。入れかえると値が変わる（8-1） \\
$\dfrac{\partial f}{\partial x}$ & $\dfrac{\dd f}{\dd x}$ & 変数が2つ以上なら $\partial$（ほかを止める）。止めた文字は数と同じで、微分しない（8-2） \\
$\dfrac{\partial^2f}{\partial x^2}$ & $\Bigl(\dfrac{\partial f}{\partial x}\Bigr)^2$ & 「2回偏微分する」の印。1回めの答えを2乗するのではない（8-2） \\
$\dd f$ & $\dfrac{\partial f}{\partial x}+\dfrac{\partial f}{\partial y}$ & 全微分は「坂×進んだ長さ」をたす。$\dd x$・$\dd y$ を忘れない（8-3） \\
$\mathrm{grad}\,h$ & 1つの数 & 矢印 $(-0.6,\ -0.8)$。2つの成分をたさない。ボールは $-\mathrm{grad}\,h$ の向きに転がる（8-4） \\
$\nabla f$ & $f\nabla$ & $\nabla$ は右にある関数を微分する。左に置くと、何も微分していない（8-4） \\
$r\,\dd r\,\dd\theta$ & $\dd r\,\dd\theta$ & 極座標のタイルの $r$ を忘れない。$\dd r\,\dd\theta$ の単位は m で、面積にならない（8-3・8-6） \\
$r^2\sin\theta$ & $r^2$ & 球座標のタイル。$\theta$ は0から $\pi$、$\varphi$ は0から $2\pi$（8-6・8-8） \\
$\displaystyle\int_0^1\!\!\int_0^2f\,\dd x\,\dd y$ & 範囲の組 & 内側の ∫ と内側の $\dd x$ が組。内側の範囲には、外側の文字だけが入る（8-6） \\
$\displaystyle\iint_D$ & 2つの積分のかけ算 & かけ算に分かれるのは、長方形で $f=g(x)\,h(y)$ のとき（場合 A$'$）だけ（8-6） \\
$\displaystyle\oint_C$ & 円の道 & ◯は「閉じた道」という印。正方形でも三角形でもよい。どの道かは $C$ で決める（8-7） \\
$\displaystyle\oiint_S$ & 球の面 & ◯は「閉じた面」という印。箱でもよい。どの面かは $S$ で決める（8-8・8-9・8-11） \\
$\displaystyle\int_C\vec F\cdot\dd\vec r$ & $|\vec F|\times$長さ & 内積をたす。力と進む向きが反対なら負、直角なら0（8-7） \\
$\displaystyle\int_0^1\dots\dd t$ & $\displaystyle\int_1^0$ & 下のはしが「はじめの点」、上のはしが「おわりの点」。逆にすると符号が逆（8-7） \\
$\vec F\cdot\vec n$ & $|\vec F|$ & 面に直角な成分だけが通りぬける。閉じた面の $\vec n$ は外向き（8-8） \\
$\mathrm{div}\,\vec v$ & 矢印 & 答えは数。$v_x$ は $x$ で、$v_y$ は $y$ で偏微分してたす（8-9） \\
$\mathrm{rot}\,\vec F$ & $\mathrm{div}\,\vec F$ & 平面は $\dfrac{\partial F_y}{\partial x}-\dfrac{\partial F_x}{\partial y}$（ひく順に注意）。3次元の回転は矢印（8-10） \\
$\nabla^2f$ & $\Delta f$、$(\nabla f)^2$ & 2回微分の和で、答えは数。この本の $\Delta$ は「変化の量」（8-12） \\
$T_i^{\,n}$ & $T_i$ の $n$ 乗 & 右上の $n$ は「$n$ 歩め」の名札（8-13） \\
$\dfrac{\alpha\,\Delta t}{\Delta x^2}$ & $\dfrac{\alpha\,\Delta t}{\Delta x}$ & $r$ の $\Delta x$ は2乗。$r\leqq\frac12$ をこえると計算がこわれる（8-13） \\
$\dfrac{\mathrm DT}{\mathrm Dt}$ & $\dfrac{\partial T}{\partial t}$ & 立てた D は、流れといっしょに動いて見た変化（物質微分）（8-13） \\
\end{longtblr}\par}

%% 記号カードのつづき：ギリシャ文字（この巻での意味に合わせた）
\par\noindent\begin{minipage}{\linewidth}
\bkbkhead{ギリシャ文字の読み方}
{\small この巻では、密度 $\rho$、熱拡散率 $\alpha$、真空の誘電率 $\varepsilon_0$ など、ギリシャ文字がたくさん出てきます。大文字・小文字・読み方と、この巻での意味です（かっこの中は章）。}\par\vspace{2mm}
\begin{bkcenter}
\footnotesize
\begin{tblr}{colspec={Q[c,m,wd=5mm] Q[c,m,wd=5mm] Q[l,m,wd=17mm] Q[l,m,wd=35mm] Q[c,m,wd=5mm] Q[c,m,wd=5mm] Q[l,m,wd=17mm] Q[l,m,wd=35mm]}, row{1}={bg=bkVol!14, font=\sffamily\bfseries\footnotesize}, hlines={0.4pt,bkLine}, vline{5}={0.8pt,bkVol}, rowsep=0.6pt, colsep=2pt}
大 & 小 & 読み方 & この巻での意味 & 大 & 小 & 読み方 & この巻での意味 \\
$\mathrm{A}$ & $\alpha$ & アルファ & 熱拡散率（8-13）、円すいの半分の角（8-8） & $\mathrm{N}$ & $\nu$ & ニュー & ― \\
$\mathrm{B}$ & $\beta$ & ベータ & フィンの式の $\beta$（8-13） & $\Xi$ & $\xi$ & クシー & ― \\
$\Gamma$ & $\gamma$ & ガンマ & ― & $\mathrm{O}$ & $o$ & オミクロン & ― \\
$\Delta$ & $\delta$ & デルタ & $\Delta x$、$\Delta t$（8-13）、$\delta Q$（8-3） & $\Pi$ & $\pi$ & パイ & 円周率 \\
$\mathrm{E}$ & $\varepsilon$ & イプシロン & 真空の誘電率 $\varepsilon_0$（8-8） & $\mathrm{P}$ & $\rho$ & ロー & 密度、電荷の密度（8-9） \\
$\mathrm{Z}$ & $\zeta$ & ゼータ & ― & $\Sigma$ & $\sigma$ & シグマ & $\Sigma$ は和、$\sigma$ は面密度（8-6） \\
$\mathrm{H}$ & $\eta$ & イータ & ― & $\mathrm{T}$ & $\tau$ & タウ & ― \\
$\Theta$ & $\theta$ & シータ & 極座標・球座標の角（真上から） & $\Upsilon$ & $\upsilon$ & ウプシロン & ― \\
$\mathrm{I}$ & $\iota$ & イオタ & ― & $\Phi$ & $\phi$ & ファイ & $\phi$ は電位（8-4）、$\Phi$ は磁束（8-11）。球座標の角 $\varphi$ は別の文字 \\
$\mathrm{K}$ & $\kappa$ & カッパ & 専門書の熱拡散率（8-13） & $\mathrm{X}$ & $\chi$ & カイ & ― \\
$\Lambda$ & $\lambda$ & ラムダ & 固有値（8-13）、線密度（8-7）、線電荷密度（8-11） & $\Psi$ & $\psi$ & プサイ & ― \\
$\mathrm{M}$ & $\mu$ & ミュー & 真空の透磁率 $\mu_0$（8-10）、動摩擦係数 $\mu'$（8-7）、粘度（8-14） & $\Omega$ & $\omega$ & オメガ & $\omega$ は角速度（8-10）、$\Omega$ は立体角（8-8） \\
\end{tblr}
\end{bkcenter}
{\footnotesize 大文字が英語のアルファベットと同じ形の文字（A、B、E など）は、記号としてはほとんど使いません。電位の $\phi$ と球座標の角 $\varphi$ は別の文字です（同じページに出るときは、形で見分けます）。$\nabla$（ナブラ）はギリシャ文字ではなく、たて琴の形からついた名前です（8-4）。熱伝導率は $k$、熱拡散率は $\alpha$ と書きます（$\lambda$ は固有値とぶつかるので使いません）。}
\end{minipage}




%% ---------------- 定数・単位・公式 ----------------
%% 積分の地図は見開き（左右2ページ）なので、見出しのページを偶数ページ（左）にそろえる
\bkevenpage
\bkbacktitle{よく使う定数・単位・公式}
\bksubsection{1}{積分の地図}
\label{bk:8-back:chizu}
{\kviiizbody\raggedright 積分の記号は「タイルに分けてたし、細かくした先」という操作の\textbf{名前}です。この見開きに、すべての場合の「こんなとき→直す手順→例題（その表の章の番号）」をまとめました。\par}
\vspace{1mm}

%% 見分け方（名前と、ふつうの定積分。◯は閉じている印）と、どの表を見る？（流れ図）
\noindent\begin{minipage}{\linewidth}
\kviiizhead{名前と、ふつうの定積分の見分け方}
\begin{center}
\kviiizNt{\iint_D\sigma\,\dd A}\ $=$\ \kviiizSt{\int_0^1\!\Bigl(\int_0^2(10+5x+2y)\,\dd x\Bigr)\dd y}\ $=32$\quad{\footnotesize（8-6 のはてな？の板。④）}
\end{center}
\vspace{-1.5mm}
\noindent\begin{minipage}[t]{0.41\linewidth}\raggedright\kviiizbody\vspace{0pt}
{\sffamily\bfseries 点線の枠＝名前}\par
まだ計算できない記号。何を（$\sigma$）・どこで（$D$）・1枚の量（$\dd A$）を読む。\par\vspace{0.6mm}
{\sffamily\bfseries 実線の枠＝ふつうの定積分}\par
数の範囲がついた ∫ のくり返し（第5巻の $\int_a^b$）。ここではじめて計算できる。\par\vspace{0.6mm}
{\sffamily\bfseries ◯は「閉じている」印}\par
$\oint_C$ は閉じた道、$\oiint_S$ は閉じた面。形（円・球）や面積ではない。どの道・面かは下の $C$・$S$ で決める。\par\vspace{0.6mm}
{\sffamily\bfseries ∫ の数＝場所を表す数の数}\par
道は1つ（$t$）で $\int_C$、面は2つで $\iint$、立体は3つで $\iiint$。
\end{minipage}\hfill
\begin{minipage}[t]{0.57\linewidth}\vspace{0pt}
{\sffamily\bfseries\small\color{\bkTT}どの表を見る？}\par\vspace{0.5mm}
\begin{tikzpicture}[x=1mm, y=1mm, kviiiz q/.append style={aspect=4.4}]
  \node[kviiiz do, text width=37mm] (s) at (19,0) {記号を読む（名前の段）};
  \node[kviiiz q] (q1) at (19,-12) {閉じている？};
  \node[kviiiz do, anchor=west, text width=29mm] (r1) at (42,-12) {\kviiiztab{8-11}\,で楽な側を};
  \node[kviiiz q, aspect=2.6] (q1b) at (58.5,-25) {楽になった？};
  \node[kviiiz q] (q2) at (19,-38) {面の積分？};
  \node[kviiiz do, anchor=west, text width=29mm] (r2) at (42,-38) {$\vec n\,\dd S$・$\dd S$\enspace\kviiiztab{8-8}};
  \node[kviiiz q] (q3) at (19,-51) {道の積分？};
  \node[kviiiz do, anchor=west, text width=29mm] (r3) at (42,-51) {$\dd\vec r$・$\dd s$\enspace\kviiiztab{8-7}};
  \node[kviiiz do, text width=33mm] (r4) at (19,-64) {$\dd A$・$\dd V$\enspace\kviiiztab{8-6}};
  \node[kviiiz do, anchor=west, text width=29mm] (end) at (42,-64) {確かめ：ちがう方法\\で同じ答えか};
  \draw[kviiiz go] (s) -- (q1);
  \draw[kviiiz go] (q1) -- node[kviiiz yn, above] {はい} (r1);
  \draw[kviiiz go] (r1) -- (q1b);
  \draw[kviiiz go] (q1) -- node[kviiiz yn, left=1mm] {いいえ} (q2);
  \draw[kviiiz go] (q2) -- node[kviiiz yn, above] {はい} (r2);
  \draw[kviiiz go] (q2) -- node[kviiiz yn, left=1mm] {いいえ} (q3);
  \draw[kviiiz go] (q3) -- node[kviiiz yn, above] {はい} (r3);
  \draw[kviiiz go] (q3) -- node[kviiiz yn, left=1mm] {いいえ} (r4);
  \draw[kviiiz go] (r4) -- (end);
  % 楽にならなければ、左の列の次の問いの上へもどる
  \draw[kviiiz go, shorten >=0.4mm] (q1b.south) -- node[kviiiz yn, right=0.6mm] {いいえ} (58.5,-31.5) -| (19.8,-31.5);
  % 楽になった・計算した四角は、右はしの線にまとめて「確かめ」へ
  \draw[bkInk!60, line width=0.7pt, shorten <=1.5mm] (q1b.east) -- node[kviiiz yn, above] {はい} (78,-25);
  \foreach \n in {r2,r3} {\draw[bkInk!60, line width=0.7pt, shorten <=1.5mm] (\n.east) -- (\n.east -| 78,0);}
  \draw[bkInk!60, line width=0.7pt] (78,-25) -- (78,-56.5);
  \draw[kviiiz go, shorten <=0pt] (78,-56.5) -| ([xshift=4mm]end.north);
\end{tikzpicture}
\end{minipage}\par\vspace{1mm}
\noindent{\kviiizbody どの積分も\enspace\kviiizdan{1}\enspace→\enspace\kviiizdan{3}\enspace→\enspace\kviiizdan{2}\enspace→\enspace 確かめ\enspace の順に考えます。\par}
\end{minipage}
\vspace{1.5mm}

%% 表8-6 重積分
\noindent\begin{kviiizmap}{row{2}={bg=bkVol!14, font=\sffamily\bfseries\footnotesize}}
\SetCell[c=3]{l} \kviiizband{\kviiiztab{8-6}\enspace{\sffamily\bfseries\color{\bkTT}重積分}\enspace\kviiizpic{2}\,\kviiizN{\iint_D f\,\dd A}\enspace\kviiizpic{3}\,\kviiizN{\iiint_V f\,\dd V}\enspace\kviiizto\enspace 定積分を2回（立体は3回）} & & \\
記号 & こんなとき\enspace→\enspace ふつうの定積分に直す手順 & 例題 \\
\kviiizcase{A} & \kviiizif{長方形}$\int_c^d\bigl(\int_a^b f\,\dd x\bigr)\dd y$。内側は $y$ を止める。順番は自由 & 1 \\
\kviiizcase{A$'$} & \kviiizif{長方形で、$f$ が「$x$ だけ」×「$y$ だけ」}2つの定積分のかけ算 & 2 \\
\kviiizcase{B} & \kviiizif{境目がななめの線・曲線}たて線で切り、内側の範囲を外側の文字の式で & 3・4 \\
\kviiizcase{C} & \kviiizif{円の形（円・おうぎ形・輪）}極座標 $\dd A=r\,\dd r\,\dd\theta$ で2回 & 5 \\
\kviiizcase{C$'$} & \kviiizif{円の形で、$f$ は中心からの距離 $r$ だけ}細い輪 $2\pi r\,\dd r$ で1回 & 6・7 \\
\kviiizcase{D} & \kviiizif{体積で、切り口の面積 $A(z)$ がわかる}体積＝$\int A(z)\,\dd z$ & 8 \\
\kviiizcase{E} & \kviiizif{立体（直方体・円柱・球など）}3回（円柱座標・球座標） & 9 \\
\kviiizcase{F} & \kviiizif{簡単な形の和・差}分けて、それぞれ A〜E で計算し、たす・ひく & 10 \\
\kviiizcase{G} & \kviiizif{折り返し・回転で0や同じになる}先に使う。左右対称な範囲で $x$ の奇関数は0 & 5・6 \\
\kviiizcase{H} & \kviiizif{データだけ}表計算で「マスの値×マスの面積」をたす & 11 \\
\end{kviiizmap}
\par
\clearpage

%% ---- 積分の地図（右のページ）：表8-7・表8-8・表8-11 ----
%% 表8-7 線積分（行の順は本文の表と同じ）
\noindent\begin{kviiizmap}{row{2}={bg=bkVol!14, font=\sffamily\bfseries\footnotesize}}
\SetCell[c=3]{l} \kviiizband{\kviiiztab{8-7}\enspace{\sffamily\bfseries\color{\bkTT}線積分}\enspace\kviiizpic{5}\,\kviiizN{\int_C\vec F\cdot\dd\vec r}\enspace\kviiizpic{4}\,\kviiizN{\int_C f\,\dd s}\enspace\kviiizto\enspace 道を $t$ で表し、定積分1回} & & \\
記号 & こんなとき\enspace→\enspace ふつうの定積分に直す手順 & 例題 \\
\kviiizcase{A} & \kviiizif{軸に平行な線分だけ}横は $F_x\,\dd x$（$y$ は止める）、たては $F_y\,\dd y$ & 1・8 \\
\kviiizcase{B} & \kviiizif{まっすぐなななめ}$\vec r=\vec A+t(\vec B-\vec A)$（$0\leqq t\leqq1$）、$\dd\vec r=(\vec B-\vec A)\,\dd t$ & 2 \\
\kviiizcase{C} & \kviiizif{$y=g(x)$ のグラフ}$x$ で：$\dd\vec r=(1,\ g'(x))\,\dd x$。$y$ に $g(x)$ を入れる & 3 \\
\kviiizcase{D} & \kviiizif{円・円弧}角 $\theta$ で：$\dd\vec r=(-R\sin\theta,\ R\cos\theta)\,\dd\theta$ & 4・9 \\
\kviiizcase{E} & \kviiizif{ななめ・曲線のある折れ線}区間ごとに A〜D でたす。逆向きは符号が逆 & 5 \\
\kviiizcase{I} & \kviiizif{長さでたす $\int_C f\,\dd s$}$\dd s=|\vec r\,'(t)|\,\dd t$。向きを逆にしても同じ値 & 9 \\
\kviiizcase{F} & \kviiizif{力が式で書いてない}直角なら0、平行で一定なら $\pm F\times$（長さ） & 6 \\
\kviiizcase{G} & \kviiizif{$U$ のわかる保存力}道によらない：$U(\text{はじめ})-U(\text{おわり})$ & 7 \\
\kviiizcase{H} & \kviiizif{閉じた道（$\oint$）}保存力なら0。ちがえば1周を F〜D の問いで & 4・5・8 \\
\kviiizcase{J} & \kviiizif{データだけ・数値で}表計算で $\sum\vec F_i\cdot\Delta\vec r_i$（まん中の点で） & 10 \\
\kviiizcase{K} & \kviiizif{$p$-$V$ 図の上の道}区間ごとに $\int p\,\dd V$（直し方は A と同じ） & 11 \\
\end{kviiizmap}
\par\vspace{2mm}

%% 表8-8 面積分
\noindent\begin{kviiizmap}{row{2}={bg=bkVol!14, font=\sffamily\bfseries\footnotesize}}
\SetCell[c=3]{l} \kviiizband{\kviiiztab{8-8}\enspace{\sffamily\bfseries\color{\bkTT}面積分}\enspace\kviiizpic{7}\,\kviiizN{\iint_S\vec F\cdot\vec n\,\dd S}\enspace\kviiizpic{6}\,\kviiizN{\iint_S f\,\dd S}\enspace\kviiizto\enspace 面を2つの数で表し、重積分に} & & \\
記号 & こんなとき\enspace→\enspace ふつうの定積分に直す手順 & 例題 \\
\kviiizcase{A} & \kviiizif{平らな面で、場が一様}積分しない：$\vec F\cdot(A\vec n)$（$A\vec n$ は外積で） & 1・2・3 \\
\kviiizcase{B} & \kviiizif{平らな面（円でない）}$\vec n$ は軸の向き。8-6 の重積分そのもの & 4 \\
\kviiizcase{C} & \kviiizif{平らな円の面}極座標 $r\,\dd r\,\dd\theta$。$r$ だけなら $2\pi r\,\dd r$ & 5 \\
\kviiizcase{D} & \kviiizif{球面・円柱面で、$\vec F\cdot\vec n$ が一定}一定の値×面積（$4\pi R^2$、$2\pi RL$） & 6・9 \\
\kviiizcase{E} & \kviiizif{$z=g(x,\ y)$ のグラフ}$\vec n\,\dd S=(-\partial g/\partial x,\ -\partial g/\partial y,\ 1)\,\dd x\,\dd y$（上向き） & 7 \\
\kviiizcase{F} & \kviiizif{球面・円柱面}$(\vec r_u\times\vec r_v)\,\dd u\,\dd v$。球面 $\dd S=R^2\sin\theta\,\dd\theta\,\dd\varphi$ & 8・9 \\
\kviiizcase{G} & \kviiizif{閉じた面で、箱の形}面ごとに B でたす。どの面も外向きが正 & 10 \\
\kviiizcase{H} & \kviiizif{閉じた面で、点電荷1つの場}中なら同じ中心の球（D）、外なら0 & 11 \\
\kviiizcase{I} & \kviiizif{立体角を求める}$\Omega=\iint_S(\vec e_r\cdot\vec n)/r^2\,\dd S$。円すいは $2\pi(1-\cos\alpha)$ & 12 \\
\kviiizcase{J} & \kviiizif{データだけ}表計算で「$\vec F\cdot\vec n$×タイルの面積」をたす & 13 \\
\end{kviiizmap}
\par\vspace{2mm}

%% 表8-11 ガウスの定理とストークスの定理
\noindent\begin{kviiizmap}{row{2}={bg=bkVol!14, font=\sffamily\bfseries\footnotesize}, row{10}={bg=white}}
\SetCell[c=3]{l} \kviiizband{\kviiiztab{8-11}\enspace{\sffamily\bfseries\color{\bkTT}ガウスの定理}\enspace\kviiizN{\oiint_S\vec F\cdot\vec n\,\dd S}\,$=$\,\kviiizN{\iiint_V\mathrm{div}\,\vec F\,\dd V}\enspace\kviiizto\enspace 楽な側を選び、表8-6・8-8 で直す} & & \\
記号 & こんなとき\enspace→\enspace 楽な側に取りかえて、直す & 例題 \\
\kviiizcase{G1} & \kviiizif{div は簡単}体積の積分に。div が定数なら「div×体積」 & 1・2 \\
\kviiizcase{G2} & \kviiizif{開いた面}ふたをして G1。ふたの分（8-8 の B）をひく & 3 \\
\kviiizcase{G3} & \kviiizif{対称な場の大きさを知りたい}球・円柱・箱の袋で、大きさ×面積＝中の量 & 4・5 \\
\kviiizcase{G4} & \kviiizif{袋の一部の面の流束}全体は中の量で決まる。同じ形の面に等分 & 6 \\
\kviiizcase{G5} & \kviiizif{div が決まらない点がある}その点を小さな球（平面なら円）でくりぬく & 7 \\
\kviiizcase{G6} & \kviiizif{平面（2次元）の流れ}$\oint_C\vec F\cdot\vec n\,\dd s=\iint_D\mathrm{div}\,\vec F\,\dd A$ & 8 \\
\kviiizcase{G7} & \kviiizif{圧力のような数の面の積分}$\oiint_S p\,\vec n\,\dd S=\iiint_V\mathrm{grad}\,p\,\dd V$ & 9 \\
\SetCell[c=3]{l} \kviiizband{\kviiiztab{8-11}\enspace{\sffamily\bfseries\color{\bkTT}ストークスの定理}\enspace\kviiizN{\oint_C\vec F\cdot\dd\vec r}\,$=$\,\kviiizN{\iint_S\mathrm{rot}\,\vec F\cdot\vec n\,\dd S}\enspace\kviiizto\enspace 道を回る向きと $\vec n$ は右ねじ} & & \\
\kviiizcase{S1} & \kviiizif{rot は簡単}面積分に。一定なら「rot×面積」 & 10 \\
\kviiizcase{S2} & \kviiizif{面積を知りたい}$\oint_C(x\,\dd y-y\,\dd x)$ の半分。多角形は、くつひもの式 & 11 \\
\kviiizcase{S3} & \kviiizif{rot＝0（穴の点をのぞく）}かこまなければ0、かこめば小さな円と同じ & 12 \\
\kviiizcase{S4} & \kviiizif{軸対称な場の大きさを知りたい}同じ軸の円を道に：$B\times2\pi r=\mu_0I$ & 13 \\
\kviiizcase{S5} & \kviiizif{面の積分がほしい（ふちは同じ）}計算が楽な面（平らな円板など）を選ぶ & 14 \\
\kviiizcase{S6} & \kviiizif{時間で変わる磁場がある}$\oint_C\vec E\cdot\dd\vec r=-\dd\Phi/\dd t$（読むだけ） & 15 \\
\end{kviiizmap}
\par
\clearpage


\bksubsection{2}{この巻で使う単位}
{\small
\begin{longtblr}{colspec={Q[l,m,wd=26mm] Q[l,m,wd=30mm] Q[l,m,wd=53mm] Q[c,m,wd=12mm]}, rowhead=1, row{1}={bg=bkVol!14, font=\sffamily\bfseries\footnotesize}, hlines={0.4pt,bkLine}, rowsep=1.4pt}
量 & 単位（読み方） & ほかの単位との関係・この巻の例 & 章 \\
長さ・高さ & m、cm、mm、km & 主役の山の $x$・$y$・$h$ はどれも m & 8-1 \\
面積・体積 & m²、m³、L、mL & $\qty{1}{L}=\qty{0.001}{m^3}$、\ $\qty{1}{mL}=\qty{1}{cm^3}$ & 8-6 \\
偏微分 & （出てくる量の単位）÷（動かした量の単位） & 坂 $\dfrac{\partial h}{\partial x}$ は m/m、$\dfrac{\partial T}{\partial x}$ は K/m、$\dfrac{\partial p}{\partial z}$ は Pa/m & 8-2 \\
面密度・線密度 & kg/m²、kg/m & 板 $\sigma=10+5x+2y$〔kg/m²〕、針金 $\lambda=2+\sin\theta$〔kg/m〕 & {8-6\\8-7} \\
仕事・エネルギー & J（ジュール） & $\qty{1}{J}=\qty{1}{N.m}$。力の線積分の単位 & 8-7 \\
流量 & m³/s、mL/s & 管の流量 \qty{31.4}{mL/s}（体積の流束） & {8-6\\8-8} \\
発散・回転 & 1/s & 速さ〔m/s〕を長さ〔m〕でわる。循環は m²/s & {8-9\\8-10} \\
角速度 & rad/s & 流れに置いた風車は $\dfrac12\,\mathrm{rot}\,\vec v$ & 8-10 \\
立体角 & sr（ステラジアン） & 全体は $4\pi$ sr。半分の角30°の円すいは 0.842 sr & 8-8 \\
電場 & N/C（＝V/m） & 点電荷の電場は $\dfrac{1}{r^2}$ に比例 & {8-5\\8-11} \\
電場の流束 $\Phi_E$ & N·m²/C & {$\dfrac{Q}{\varepsilon_0}$。\qty{1.0e-9}{C} で約113。電気工学の本の「電束」は、電束密度 $\vec D=\varepsilon_0\vec E$ の流束で、単位は C} & 8-8 \\
磁場（磁束密度） & T（テスラ） & 10\,A・5\,cm で $4\times10^{-5}$\,T & {8-10\\8-11} \\
熱の流れ & W/m² & 6-3 のかべ \qty{160}{W/m^2} & 8-4 \\
熱伝導率 & W/(m·K) & コンクリート 1.6（ばね定数の $k$ とは別） & 8-4 \\
熱拡散率 & m²/s & $\alpha=\dfrac{k}{\rho c}$。アルミ $1.0\times10^{-4}$ & 8-13 \\
温度・温度の差 & ℃、K（ケルビン） & 温度の差は \qty{1}{K} と 1\,℃ が同じ & {8-4\\8-13} \\
圧力 & Pa、hPa、kPa & $\qty{1}{hPa}=\qty{100}{Pa}$、\ $\qty{1}{kPa}=\qty{1000}{Pa}$ & {8-2\\8-4} \\
\end{longtblr}\par}

\bksubsection{3}{よく使う数}
{\small
\begin{longtblr}{colspec={Q[l,m,wd=40mm] Q[l,m,wd=72mm] Q[c,m,wd=12mm]}, rowhead=1, row{1}={bg=bkVol!14, font=\sffamily\bfseries\footnotesize}, hlines={0.4pt,bkLine}, rowsep=1pt}
名前 & 値 & 章 \\
$\sqrt\pi$ & 1.7725（ガウス積分 $\int_{-\infty}^{\infty}e^{-x^2}\,\dd x$） & 8-6 \\
$\dfrac{1}{\sqrt2}$ & 0.707（北東の向きの長さ1の矢印 $(0.707,\ 0.707)$） & 8-4 \\
重力加速度 $g$ & 計算では \qty{9.8}{m/s^2}。地面での値は $GM/R^2=\qty{9.82}{m/s^2}$ & {8-5\\8-11} \\
地球 & 半径 $R=\num{6371}$\,km、$GM=\qty{3.986e14}{m^3/s^2}$ & {8-5\\8-11} \\
万有引力定数 $G$ & $6.674\times10^{-11}$\,N·m²/kg² & 8-14 \\
真空の誘電率 $\varepsilon_0$ & $8.85\times10^{-12}$\,F/m。\ $\dfrac{1}{4\pi\varepsilon_0}=9.0\times10^9$\,N·m²/C² & {8-8\\8-11} \\
真空の透磁率 $\mu_0$ & {$\mu_0\approx4\pi\times10^{-7}$\,N/A²（＝T·m/A）。2019年の SI の改定からは、測って決める値で、ほぼこの値} & {8-10\\8-11} \\
光の速さ $c$ & $\dfrac{1}{\sqrt{\mu_0\varepsilon_0}}$、約 \qty{3.00e8}{m/s} & 8-14 \\
地球の磁場（日本） & およそ $4.5\times10^{-5}$\,T & 8-10 \\
水の密度 & \qty{1000}{kg/m^3} & {8-2\\8-11} \\
熱拡散率 $\alpha$〔m²/s〕 & {銅 $1.2\times10^{-4}$、アルミ $1.0\times10^{-4}$、鉄 $2.3\times10^{-5}$、ステンレス $4\times10^{-6}$、水（動かないとき）$1.4\times10^{-7}$。\qty{10}{cm} に広がる時間のめやす $L^2/\alpha$ は、順に約80秒・100秒・7分・40分・20時間} & 8-13 \\
\end{longtblr}\par}

\Needspace{8\baselineskip}
\bkbkhead{この巻でくり返し使う場面と数}
{\small 同じ場面は、どの章でも同じ数で計算しています。}\par\vspace{1mm}
{\small
\begin{longtblr}{colspec={Q[l,m,wd=21mm] Q[l,m,wd=92mm] Q[c,m,wd=12mm]}, rowhead=1, row{1}={bg=bkVol!14, font=\sffamily\bfseries\footnotesize}, hlines={0.4pt,bkLine}, rowsep=1.6pt}
場面 & 式と数 & 章 \\
主役の山 & $h=100-0.01x^2-0.02y^2$〔m〕（$x$ は東へ、$y$ は北へ）。点 $(30,\ 20)$ で $h=83$、$\dfrac{\partial h}{\partial x}=-0.6$、$\dfrac{\partial h}{\partial y}=-0.8$、$\mathrm{grad}\,h=(-0.6,\ -0.8)$（大きさ1.0）。東へも北へも \qty{1}{m} で全微分 $-1.4$（本当は $-1.43$）。体積 $1.11\times10^6$\,m³。\qty{50}{kg} の人が頂上まで登る仕事 \qty{8330}{J} & {8-1\\〜8-4\\8-6\\8-7} \\
2つの力 & $\vec F_1=(y,\ x)$、$\vec F_2=(-y,\ x)$（物理第1巻と同じ）。道A $(0,\ 0)\to(1,\ 0)\to(1,\ 1)$、道B $(0,\ 0)\to(0,\ 1)\to(1,\ 1)$。$\vec F_1$ はどちらも1（保存力、$U=-xy$）。$\vec F_2$ は道A 1・道B $-1$、正方形の1周2、半径1の円の1周 $2\pi$（rot は2） & {8-7\\8-10\\8-11} \\
池 & わき出しの口 $(1.5,\ 2.5)$、吸いこみの口 $(4.5,\ 1.5)$、強さ $Q=4$\,m²/s。口のマスの発散は4と $-4$、ほかのマスは0。口をかこむどんな道からも4が出る & {8-5\\8-9\\8-11} \\
ロボットの腕 & $L_1=\qty{0.30}{m}$、$L_2=\qty{0.25}{m}$、30°・45°。$J=\begin{pmatrix}-0.3915&-0.2415\\0.3245&0.0647\end{pmatrix}$、$\det J=0.0530$ & 8-3 \\
厚さがちがう板 & $\sigma=10+5x+2y$〔kg/m²〕、\qty{2}{m}×\qty{1}{m} の板で \qty{32}{kg}。4×2 のタイル（まん中の点）でも32、左下の角の値では28.5 & 8-6 \\
管の流れ & 半径 $R=\qty{1}{cm}$、$u=0.2\Bigl(1-\dfrac{r^2}{R^2}\Bigr)$〔m/s〕。流量 \qty{31.4}{mL/s}、平均の速さは最大の半分 & {8-6\\8-8} \\
6-3 のかべ & 厚さ \qty{0.2}{m}、20\,℃と0\,℃、$k=1.6$。$T=20-100x$、$\vec q=-k\,\mathrm{grad}\,T=(160,\ 0,\ 0)$〔W/m²〕 & 8-4 \\
立方体の箱 & 1辺1、$\vec F=(x^2,\ xy,\ -z)$。6つの面から外へ出る量の合計0.5。$\mathrm{div}\,\vec F=3x-1$ の体積の積分も0.5 & {8-8\\8-11} \\
点電荷と地球 & 電場の流束 $\Phi_E=\dfrac{Q}{\varepsilon_0}$（袋の大きさによらない）。地球の重力は外では全部の質量が中心にあるのと同じ。一様な球なら、半分の深さで \qty{4.91}{m/s^2} & {8-8\\8-11} \\
直線電流 & 10\,A、\qty{5}{cm} の所で $B=\dfrac{\mu_0I}{2\pi r}=4\times10^{-5}$\,T（地球の磁場と同じくらい） & {8-10\\8-11} \\
ならす表計算 & 5×5 の点、上のふち100、ほかのふち0。くり返すと、まん中は25 & 8-12 \\
棒の熱の表計算 & $\alpha=1\times10^{-4}$\,m²/s、$\Delta x=\qty{0.01}{m}$、$\Delta t=\qty{0.25}{s}$ で $r=0.25$。左はし100\,℃、はじめ20\,℃、右はしは熱が出ない。まん中が50\,℃をこえるのは15.5秒後 & 8-13 \\
\end{longtblr}\par}






\noindent\begin{minipage}{\linewidth}
\noindent{\sffamily\bfseries\small\color{\bkTT}主役の6つの場のわき出しとうず}（8-5・8-9・8-10・8-14）\par\vspace{1mm}
\begin{bkcenter}
\begin{tblr}{colspec={Q[l,m,wd=28mm] *{6}{Q[c,m,wd=16mm]}}, row{1-2}={bg=bkVol!14, font=\sffamily\bfseries\footnotesize}, hlines={0.4pt,bkLine}, rowsep=1.2pt, colsep=2.5pt, cells={font=\small}}
場 & （ア） & （イ） & （ウ） & （エ） & （オ） & （カ） \\
 & 広がる & 集まる & そろった & うず & ずれ & 側溝 \\
矢印 $\vec F$ & $(x,\ y)$ & $(-x,\ -y)$ & $(1,\ 0)$ & $(-y,\ x)$ & $(y,\ 0)$ & $(x,\ 0)$ \\
わき出し $\mathrm{div}\,\vec F$ & 2 & $-2$ & 0 & 0 & 0 & 1 \\
うず $\mathrm{rot}\,\vec F$ & 0 & 0 & 0 & 2 & $-1$ & 0 \\
\end{tblr}
\end{bkcenter}
{\footnotesize （ア）（イ）（カ）はわき出しだけ、（エ）（オ）はうずだけ、（ウ）はどちらもありません。回っていても div は0（エ）、まっすぐでも rot があること（オ）があります。}
\end{minipage}

\bksubsection{4}{偏微分・全微分・勾配の公式}
{\small
\begin{longtblr}{colspec={Q[l,m,wd=26mm] Q[l,m,wd=86mm] Q[c,m,wd=12mm]}, rowhead=1, row{1}={bg=bkVol!14, font=\sffamily\bfseries\footnotesize}, hlines={0.4pt,bkLine}, rowsep=1.6pt}
公式・きまり & 式 & 章 \\
2つの変数の関数 & $f(x,\ y)$。等高線は $f(x,\ y)=c$。高さのちがう等高線は交わらない。3つの変数なら等値面 & 8-1 \\
偏微分 & $\dfrac{\partial f}{\partial x}$ は $y$ を止めて $x$ だけ動かした微分。止めた文字は数と同じにあつかう & 8-2 \\
表から坂を出す & 東の坂 $\approx$（右どなり－自分）÷間かく、北の坂 $\approx$（上どなり－自分）÷間かく。両どなりの差÷（2×間かく）のほうが正確 & 8-2 \\
2回の偏微分 & $\dfrac{\partial^2f}{\partial x\,\partial y}=\dfrac{\partial^2f}{\partial y\,\partial x}$（順番によらない） & 8-2 \\
全微分 & $\dd f=\dfrac{\partial f}{\partial x}\,\dd x+\dfrac{\partial f}{\partial y}\,\dd y$（1つずつの変化をたす） & 8-3 \\
誤差の伝わり方 & $V=\pi r^2h$ なら $\dfrac{\Delta V}{V}\approx2\dfrac{\Delta r}{r}+\dfrac{\Delta h}{h}$。いちばん大きいずれは、割合の大きさをたす & 8-3 \\
連鎖律 & $\dfrac{\dd h}{\dd t}=\dfrac{\partial h}{\partial x}\dfrac{\dd x}{\dd t}+\dfrac{\partial h}{\partial y}\dfrac{\dd y}{\dd t}$（道ごとにかけて、たす） & 8-3 \\
ヤコビ行列 & 偏微分を並べた行列（行＝出てくる量、列＝動かす量）。$\dd\vec p=J\,\dd\vec\theta$ & 8-3 \\
ヤコビアン & $\det J$＝小さな面積の倍率。極座標は $r$、だ円の座標は $ab\,r$ & {8-3\\8-6} \\
勾配 & $\mathrm{grad}\,f=\Bigl(\dfrac{\partial f}{\partial x},\ \dfrac{\partial f}{\partial y},\ \dfrac{\partial f}{\partial z}\Bigr)=\nabla f$ & 8-4 \\
方向微分 & $\mathrm{grad}\,f\cdot\vec e=|\mathrm{grad}\,f|\cos\theta$（$\vec e$ は長さ1の矢印） & 8-4 \\
いちばん急な向き & 上りは grad、下りは $-\mathrm{grad}$。坂の大きさは $|\mathrm{grad}\,f|$。grad は等高線・等値面と直角 & 8-4 \\
物理の式 & $\vec F=-\mathrm{grad}\,U$、\ $\vec E=-\mathrm{grad}\,\phi$、\ $\vec q=-k\,\mathrm{grad}\,T$ & 8-4 \\
ナブラ & $\nabla=\Bigl(\dfrac{\partial}{\partial x},\ \dfrac{\partial}{\partial y},\ \dfrac{\partial}{\partial z}\Bigr)$。右にある関数を偏微分して並べる「機械」 & 8-4 \\
\end{longtblr}\par}

\bksubsection{5}{ベクトル場・発散・回転・ラプラシアンの公式}
\begin{bkcenter}
\small
\begin{tblr}{colspec={Q[c,m,wd=20mm] Q[c,m,wd=17mm] Q[c,m,wd=17mm] Q[c,m,wd=17mm] Q[l,m,wd=41mm] Q[c,m,wd=10mm]}, row{1}={bg=bkVol!14, font=\sffamily\bfseries\footnotesize}, hlines={0.4pt,bkLine}, rowsep=1.2pt, colsep=3pt}
名前 & 入れるもの & 出てくるもの & $\nabla$ で & 意味 & 章 \\
$\mathrm{grad}\,f$ & 数の場 & 矢印 & $\nabla f$ & いちばん急な上り & 8-4 \\
$\mathrm{div}\,\vec F$ & 矢印の場 & 数 & $\nabla\cdot\vec F$ & わき出し（出－入）÷大きさ & 8-9 \\
$\mathrm{rot}\,\vec F$ & 矢印の場 & 矢印（平面は数） & $\nabla\times\vec F$ & うず（循環÷面積） & 8-10 \\
$\nabla^2f$ & 数の場 & 数 & $\nabla\cdot\nabla f$ & まわりの平均との差 & 8-12 \\
\end{tblr}
\end{bkcenter}
{\small
\begin{longtblr}{colspec={Q[l,m,wd=26mm] Q[l,m,wd=86mm] Q[c,m,wd=12mm]}, rowhead=1, row{1}={bg=bkVol!14, font=\sffamily\bfseries\footnotesize}, hlines={0.4pt,bkLine}, rowsep=1.6pt}
公式・きまり & 式 & 章 \\
ベクトル場 & $\vec F(x,\ y)=(F_x,\ F_y)$。場所を入れると矢印が出てくる。矢印の図は、同じ割合で縮めてかく & 8-5 \\
流線・力線 & $\dfrac{\dd y}{\dd x}=\dfrac{F_y}{F_x}$。矢印をつないだ線で、交わらない & 8-5 \\
点電荷・重力 & 大きさは $\dfrac{1}{r^2}$ に比例。電場は外向き、重力は内向き & 8-5 \\
軸を回して見る & 矢印は同じ。成分は $R^{\mathsf T}\vec F$ で変わる & 8-5 \\
発散 & $\mathrm{div}\,\vec F=\dfrac{\partial F_x}{\partial x}+\dfrac{\partial F_y}{\partial y}+\dfrac{\partial F_z}{\partial z}=\nabla\cdot\vec F$（答えは数） & 8-9 \\
発散の定義 & $\displaystyle\mathrm{div}\,\vec F=\lim_{\Delta V\to0}\frac{1}{\Delta V}\oiint_S\vec F\cdot\vec n\,\dd S$（小さな箱から出ていく量÷体積） & 8-9 \\
連続の式 & $\dfrac{\partial\rho}{\partial t}+\mathrm{div}(\rho\vec v)=0$。縮まない流れなら $\mathrm{div}\,\vec v=0$ & 8-9 \\
回転（平面） & $\mathrm{rot}\,\vec F=\dfrac{\partial F_y}{\partial x}-\dfrac{\partial F_x}{\partial y}$＝循環÷面積。正なら左回り & 8-10 \\
回転（3次元） & {$\mathrm{rot}\,\vec F=\nabla\times\vec F=\begin{pmatrix}\dfrac{\partial F_z}{\partial y}-\dfrac{\partial F_y}{\partial z}\\[8pt]\dfrac{\partial F_x}{\partial z}-\dfrac{\partial F_z}{\partial x}\\[8pt]\dfrac{\partial F_y}{\partial x}-\dfrac{\partial F_x}{\partial y}\end{pmatrix}$\\向きは回る軸。右ねじの向き} & 8-10 \\
風車 & 角速度は $\dfrac12\,\mathrm{rot}\,\vec v$。$\mathrm{rot}\,(\vec\omega\times\vec r)=2\vec\omega$ & 8-10 \\
いつも0 & $\mathrm{rot}\,(\mathrm{grad}\,f)=\vec0$、\ $\mathrm{div}\,(\mathrm{rot}\,\vec F)=0$ & 8-10 \\
保存力の見分け方 & 平面の力なら $\dfrac{\partial F_y}{\partial x}-\dfrac{\partial F_x}{\partial y}=0$（すきまのない広がりの中なら） & {8-7\\8-10} \\
ラプラシアン & $\nabla^2f=\dfrac{\partial^2f}{\partial x^2}+\dfrac{\partial^2f}{\partial y^2}+\dfrac{\partial^2f}{\partial z^2}=\mathrm{div}\,\mathrm{grad}\,f$ & 8-12 \\
まわりの平均 & （まわり4点の平均）－（自分）$\approx\dfrac{h^2}{4}\nabla^2f$（$h$ は点の間かく）。正なら温まる & 8-12 \\
落ち着いた形 & $\nabla^2T=0$（ラプラス方程式）、$\nabla^2T=-\dfrac{Q}{k}$（ポアソン方程式） & 8-12 \\
球・円柱で $r$ だけ & 球 $\dfrac{1}{r^2}\dfrac{\dd}{\dd r}\Bigl(r^2\dfrac{\dd f}{\dd r}\Bigr)$、円柱 $\dfrac{1}{r}\dfrac{\dd}{\dd r}\Bigl(r\dfrac{\dd f}{\dd r}\Bigr)$ & 8-12 \\
\end{longtblr}\par}

\bksubsection{6}{積分と定理の公式}
{\small 積分の名前を、ふつうの定積分に直すときの手順は、積分の地図（\pageref{bk:8-back:chizu}ページ）にまとめました。ここには、よく使う結果を集めます。}\par\vspace{1mm}
{\small
\begin{longtblr}{colspec={Q[l,m,wd=26mm] Q[l,m,wd=86mm] Q[c,m,wd=12mm]}, rowhead=1, row{1}={bg=bkVol!14, font=\sffamily\bfseries\footnotesize}, hlines={0.4pt,bkLine}, rowsep=1.6pt}
公式・きまり & 式 & 章 \\
タイルの面積・体積 & 長方形 $\dd A=\dd x\,\dd y$、極座標 $\dd A=r\,\dd r\,\dd\theta$、円柱座標 $\dd V=r\,\dd r\,\dd\theta\,\dd z$、球座標 $\dd V=r^2\sin\theta\,\dd r\,\dd\theta\,\dd\varphi$ & 8-6 \\
一般の座標 & $\dd A=|\det J|\,\dd u\,\dd v$（$J$ は $(x,\ y)$ を $(u,\ v)$ で偏微分した行列） & 8-6 \\
円・球 & 円の面積 $\pi R^2$、球の体積 $\dfrac43\pi R^3$、球面の面積 $4\pi R^2$ & {8-6\\8-8} \\
断面二次モーメント & 円 $\dfrac{\pi d^4}{64}$、中空の丸棒 $\dfrac{\pi(D^4-d^4)}{64}$ & 8-6 \\
慣性モーメント & 球（中心を通る軸） $\dfrac25MR^2$ & 8-6 \\
ガウス積分 & $\displaystyle\int_{-\infty}^{\infty}e^{-x^2}\,\dd x=\sqrt\pi$（$\iint e^{-(x^2+y^2)}\,\dd A=\pi$ から） & 8-6 \\
線積分 & $\displaystyle\int_C\vec F\cdot\dd\vec r=\int_{t_1}^{t_2}\vec F(\vec r(t))\cdot\vec r\,'(t)\,\dd t$。長さでは $\dd s=|\vec r\,'(t)|\,\dd t$ & 8-7 \\
保存力 & $\vec F=-\mathrm{grad}\,U$ なら $\displaystyle\int_C\vec F\cdot\dd\vec r=U(\text{はじめ})-U(\text{おわり})$。閉じた道の1周は0 & 8-7 \\
流束 & $\displaystyle\iint_S\vec F\cdot\vec n\,\dd S$（$\vec n$ は長さ1の法線。閉じた面では外向き） & 8-8 \\
曲がった面 & $\vec n\,\dd S=(\vec r_u\times\vec r_v)\,\dd u\,\dd v$。球面 $\dd S=R^2\sin\theta\,\dd\theta\,\dd\varphi$、円柱面 $\dd S=R\,\dd\theta\,\dd z$ & 8-8 \\
立体角 & $\displaystyle\Omega=\iint_S\frac{\vec e_r\cdot\vec n}{r^2}\,\dd S$。円すい（半分の角 $\alpha$）は $2\pi(1-\cos\alpha)$、全体は $4\pi$ & 8-8 \\
ガウスの定理 & $\displaystyle\oiint_S\vec F\cdot\vec n\,\dd S=\iiint_V\mathrm{div}\,\vec F\,\dd V$（外へ出る量＝中のわき出しの合計） & 8-11 \\
平面のガウスの定理 & $\displaystyle\oint_C\vec F\cdot\vec n\,\dd s=\iint_D\mathrm{div}\,\vec F\,\dd A$（$\vec n$ は外向き） & 8-11 \\
ストークスの定理 & $\displaystyle\oint_C\vec F\cdot\dd\vec r=\iint_S\mathrm{rot}\,\vec F\cdot\vec n\,\dd S$（$C$ を回る向きと $\vec n$ は右ねじ） & 8-11 \\
グリーンの定理 & $\displaystyle\oint_C(F_x\,\dd x+F_y\,\dd y)=\iint_D\Bigl(\frac{\partial F_y}{\partial x}-\frac{\partial F_x}{\partial y}\Bigr)\dd A$（左回り） & 8-11 \\
ふちの道で面積 & $A=\dfrac12\displaystyle\oint_C(x\,\dd y-y\,\dd x)$ & 8-11 \\
ガウスの法則 & $\displaystyle\oiint_S\vec E\cdot\vec n\,\dd S=\frac{Q_{\text{中}}}{\varepsilon_0}$。点電荷 $E=\dfrac{Q}{4\pi\varepsilon_0r^2}$、まっすぐな線の電荷 $E=\dfrac{\lambda}{2\pi\varepsilon_0r}$ & 8-11 \\
アンペールの法則 & $\displaystyle\oint_C\vec B\cdot\dd\vec r=\mu_0I_{\text{中}}$。まっすぐな電流 $B=\dfrac{\mu_0I}{2\pi r}$ & 8-11 \\
浮力 & $-\displaystyle\oiint_S p\,\vec n\,\dd S$ は上向きで、大きさ $\rho gV$（物がおしのけた水の重さ） & 8-11 \\
\end{longtblr}\par}

\bksubsection{7}{偏微分方程式と、場の式}
{\small
\begin{longtblr}{colspec={Q[l,m,wd=26mm] Q[l,m,wd=86mm] Q[c,m,wd=12mm]}, rowhead=1, row{1}={bg=bkVol!14, font=\sffamily\bfseries\footnotesize}, hlines={0.4pt,bkLine}, rowsep=1.6pt}
公式・きまり & 式 & 章 \\
熱方程式 & $\dfrac{\partial T}{\partial t}=\alpha\dfrac{\partial^2T}{\partial x^2}$（3次元は $\alpha\nabla^2T$）、\ $\alpha=\dfrac{k}{\rho c}$ & 8-13 \\
1歩の式 & $T_i^{\,n+1}=T_i^{\,n}+r\bigl(T_{i+1}^{\,n}-2T_i^{\,n}+T_{i-1}^{\,n}\bigr)$、\ $r=\dfrac{\alpha\,\Delta t}{\Delta x^2}$ & 8-13 \\
こわれない条件 & $r\leqq\dfrac12$（1歩の行列の固有値の大きさが1をこえない） & 8-13 \\
$\sin$ の形の答え & $T=A\,e^{-\alpha a^2t}\sin ax$。細かいでこぼこほど、$a^2$ に比例して速く消える & 8-13 \\
広がる時間 & 長さ $L$ に熱が広がる時間のめやすは $\dfrac{L^2}{\alpha}$ & 8-13 \\
波動方程式 & $\dfrac{\partial^2y}{\partial t^2}=c^2\dfrac{\partial^2y}{\partial x^2}$。答えは $f(x-ct)+g(x+ct)$。表計算は $c\,\Delta t\leqq\Delta x$（CFL条件） & 8-13 \\
物質微分 & $\dfrac{\mathrm DT}{\mathrm Dt}=\dfrac{\partial T}{\partial t}+\vec v\cdot\mathrm{grad}\,T$（流れといっしょに動いて見た変化） & 8-13 \\
電気と磁気の4つの式 & $\mathrm{div}\,\vec E=\dfrac{\rho}{\varepsilon_0}$、\ $\mathrm{div}\,\vec B=0$、\ $\mathrm{rot}\,\vec E=-\dfrac{\partial\vec B}{\partial t}$、\ $\mathrm{rot}\,\vec B=\mu_0\vec j+\mu_0\varepsilon_0\dfrac{\partial\vec E}{\partial t}$（マクスウェル方程式） & 8-14 \\
流れの式 & $\rho\Bigl(\dfrac{\partial\vec v}{\partial t}+(\vec v\cdot\nabla)\vec v\Bigr)=-\nabla p+\mu\nabla^2\vec v+\rho\vec g$（ナビエ・ストークス方程式） & 8-14 \\
\end{longtblr}\par}

\bksubsection{8}{表計算・Python の使い方}
{\small 表計算の2次元の表は、地図と同じく北を上に置きます。1行めに $x$（東へ）、A列に $y$（北へ）を書き、行の番号が大きいほど南になります（8-1 の型）。}\par\vspace{1mm}
{\small
\begin{longtblr}{colspec={Q[l,m,wd=27mm] Q[l,m,wd=58mm] Q[l,m,wd=30mm] Q[c,m,wd=10mm]}, rowhead=1, row{1}={bg=bkVol!14, font=\sffamily\bfseries\footnotesize}, column{2}={font=\footnotesize}, hlines={0.4pt,bkLine}, rowsep=1.6pt, colsep=3pt}
すること & 書き方（例） & 意味・答え & 章 \\
{2次元の表\\（表計算）} & \texttt{=100-0.01*B\$1\^{}2-0.02*\$A2\^{}2} & {B2 に入れて、表全体へコピー} & 8-1 \\
東の差・北の差 & \texttt{=E4-D4}、\ \texttt{=D3-D4} & {右どなり－自分\\上どなり－自分} & 8-2 \\
両どなりの差 & \texttt{=(D3-B3)/20} & {左右のマスの差÷\\\qty{20}{m}（東の坂）} & 8-4 \\
タイルの和 & {\texttt{=(10+5*B\$1+2*\$A2)*0.25}\\\texttt{=SUM(B2:E3)}} & {まん中の値×タイルの面積\\合計は32} & 8-6 \\
道の線積分 & {\texttt{=-C2*D2+B2*E2}\\\texttt{=SUM(F2:F9)}} & {$\vec F_2=(-y,\ x)$ の\\$F_x\,\Delta x+F_y\,\Delta y$ をたす} & 8-7 \\
マスの循環 & \texttt{=(B2+C2-D2-E2)*1} & 下＋右－上－左 & 8-10 \\
まわりの平均との差 & \texttt{=AVERAGE(B1,A2,C2,B3)-B2} & 正なら温まる & 8-12 \\
熱の1歩 & \texttt{=C2+0.25*(B2-2*C2+D2)} & {$r=0.25$。下の行へコピー\\はしの点は境界条件で} & 8-13 \\
{場所の表を作る\\（Python）} & {\texttt{X, Y = np.meshgrid(x, y)}\\\texttt{H = 100 - 0.01*X**2 - 0.02*Y**2}} & マスごとの高さの表 & 8-1 \\
等高線の図 & \texttt{plt.contour(X, Y, H, levels=...)} & \texttt{plt.clabel} で数字 & 8-1 \\
となりとの差 & \texttt{H[:, 1:] - H[:, :-1]} & 東へ1マスの差 & 8-2 \\
矢印の図 & \texttt{plt.quiver(x, y, Fx, Fy, ...)} & 矢印の場 & 8-5 \\
タイルの和 & \texttt{(Fz * (2/n)**2).sum()} & 値×タイルの面積をたす & 8-8 \\
1歩の行列の固有値 & \texttt{np.linalg.eigvalsh(A)} & こわれないかを調べる & 8-13 \\
\end{longtblr}\par}
\begin{itemize}[leftmargin=1.4\zw, itemsep=2pt, topsep=4pt, label={\color{\bkTT}●}]
\item \texttt{\$} は、コピーしても動かない印です。\texttt{B\$1} は1行めの $x$、\texttt{\$A2} はA列の $y$ をいつも読みます。つけ忘れると、ほかのマスの値を読んでしまいます（8-1）。
\item 表計算の行の番号は下へふえますが、この本の表は北を上にしています。北の坂は「上どなり－自分」です（8-1・8-2）。
\item \texttt{np.arange(0, 60, 10)} の60は入りません（0、10、…、50）。\texttt{H[2, 3]} は、上から3行め・左から4列め（0から数える）です（8-1）。
\item タイルの和は、タイルのまん中の点の値を使うと、少ない数でも本当の値に近くなります（8-6・8-8）。\texttt{scipy} などの道具を使わず、自分でたすと、何をたしているかが見えます。
\item ふちの点（はしの点）に、まわりの平均や1歩の式を入れてはいけません。ふちは境界条件で決めます（8-12・8-13）。
\end{itemize}

\Needspace{8\baselineskip}
\bkbkhead{プログラムのひな形（8-1・8-13）}
{\small 上は主役の山の高さの表（8-1 の表計算と同じ数）。下は、棒の熱の1歩の行列（中の点9つ）の固有値（8-13）。$r=0.6$ では固有値の大きさが1をこえ、計算がこわれます。}\par\vspace{1mm}
\begin{bkcode}[Python]
import numpy as np               \# 数の表の道具 numpy を np として使う
x = np.arange(0, 60, 10)         \# 東へ 0, 10, …, 50（60 は入らない）
y = np.arange(40, -10, -10)      \# 北へ 40, 30, …, 0（北を上に）
X, Y = np.meshgrid(x, y)         \# マスごとの x の表と y の表
H = 100 - 0.01*X**2 - 0.02*Y**2  \# マスごとの高さ（B2 の式と同じ）
print(H)
\end{bkcode}
\begin{bkcodeout}[実行結果]
[[ 68.  67.  64.  59.  52.  43.]
 [ 82.  81.  78.  73.  66.  57.]
 [ 92.  91.  88.  83.  76.  67.]
 [ 98.  97.  94.  89.  82.  73.]
 [100.  99.  96.  91.  84.  75.]]
\end{bkcodeout}
\vspace{1mm}
\Needspace{66mm}
\begin{bkcode}[Python]
import numpy as np
N = 9                                \# 中の点の数
for r in [0.25, 0.5, 0.6]:
    A = (np.eye(N)*(1 - 2*r) + np.eye(N, k=1)*r
         + np.eye(N, k=-1)*r)        \# 1歩の行列
    lam = np.linalg.eigvalsh(A)      \# 固有値（小さい順）
    print(r, lam.min().round(4), lam.max().round(4))
\end{bkcode}
\begin{bkcodeout}[実行結果]
0.25 0.0245 0.9755
0.5 -0.9511 0.9511
0.6 -1.3413 0.9413
\end{bkcodeout}
\vspace{2mm}

\bksubsection{9}{専門書の書き方との対応}
{\small 大学の本では、同じものを別の記号やことばで書きます。各章の「専門書をのぞいてみよう」で読みかえたものを集めました。}\par\vspace{1mm}
{\small
\begin{longtblr}{colspec={Q[l,m,wd=52mm] Q[l,m,wd=58mm] Q[c,m,wd=14mm]}, rowhead=1, row{1}={bg=bkVol!14, font=\sffamily\bfseries\footnotesize}, hlines={0.4pt,bkLine}, rowsep=1.2pt}
専門書の書き方 & この本の書き方・意味 & 章 \\
太字の $\symbf{F}$、$\symbf{E}$ & 矢印の $\vec F$、$\vec E$ & {8-0\\8-14} \\
$\mathbb{R}^2$、等位面 & 2つの数の組（平面）全部、等値面 & 8-1 \\
$f_x$、$z_x$、$\partial_xf$ & $\dfrac{\partial f}{\partial x}$（この本では使わない） & 8-2 \\
$f_{xy}$ & $\dfrac{\partial^2f}{\partial y\,\partial x}$（$x$ で偏微分してから $y$ で） & 8-2 \\
$\Bigl(\dfrac{\partial U}{\partial T}\Bigr)_V$ & $V$ を止めて、$T$ で偏微分 & 8-3 \\
$\delta Q$、$\delta W$ & 道すじで値が変わる小さな量（全微分ではない） & 8-3 \\
$\dfrac{\partial(x,\ y)}{\partial(u,\ v)}$ & ヤコビアン $\det J$ & {8-3\\8-6} \\
$\symbf{i}$、$\symbf{j}$、$\symbf{k}$ & $\vec e_x$、$\vec e_y$、$\vec e_z$ & {8-4\\8-5} \\
$\nabla\phi$、方向微分係数 & $\mathrm{grad}\,\phi$、方向微分 & 8-4 \\
スカラー場 & 場所ごとに1つの数が決まるもの（高さ・温度・電位） & {8-4\\8-5} \\
$\dfrac{dx}{A_x}=\dfrac{dy}{A_y}=\dfrac{dz}{A_z}$ & 流線（力線）の式 & 8-5 \\
$\displaystyle\int_D f\,dA$、$\displaystyle\int_V\rho\,dV$（∫ を1つ） & $\displaystyle\iint_D f\,\dd A$、$\displaystyle\iiint_V\rho\,\dd V$ & {8-6\\8-14} \\
$\displaystyle\int_c^d dy\int_a^b f\,dx$ & $\displaystyle\int_c^d\Bigl(\int_a^b f\,\dd x\Bigr)\dd y$（右の積分から先に） & 8-6 \\
累次積分 & ふつうの定積分のくり返し & 8-6 \\
$\displaystyle\int_C P\,dx+Q\,dy$ & $\displaystyle\int_C(F_x\,\dd x+F_y\,\dd y)$ & 8-7 \\
$\symbf{F}=\nabla\phi$（数学の本の符号） & 線積分は「おわり－はじめ」。物理の本は $\vec F=-\mathrm{grad}\,U$ で「はじめ－おわり」 & 8-7 \\
$\displaystyle\oint\vec E\cdot d\vec l$、$d\vec s$ & $\displaystyle\oint_C\vec E\cdot\dd\vec r$ & 8-7 \\
閉じた面の $\displaystyle\oint_S$、$\displaystyle\oint_{\partial V}$ & $\displaystyle\oiint_S$（◯はどちらも「閉じている」印） & {8-8\\8-11} \\
$d\symbf{S}$、$d\vec A$（面素ベクトル） & $\vec n\,\dd S$ & {8-8\\8-11} \\
$\Phi_E$（電場の流束。本によっては電束） & 電場の流束。電気工学の「電束」は $\vec D$ の流束（単位 C） & 8-8 \\
$\operatorname{div}\symbf{A}\equiv\nabla\cdot\symbf{A}$ & $\mathrm{div}\,\vec A$（$\equiv$ は「こう決める」） & 8-9 \\
非圧縮性流体 & 縮まない流れる物（$\mathrm{div}\,\vec v=0$） & 8-9 \\
$\operatorname{curl}\symbf{A}$、渦度 & $\mathrm{rot}\,\vec A$、速度の場の回転 & 8-10 \\
$\partial V$、ガウスの発散定理 & $V$ の境目（ふち）、ガウスの定理 & 8-11 \\
向きと整合的 & 右ねじの関係 & 8-11 \\
$\Delta$（ラプラス作用素）、調和関数 & $\nabla^2$、$\nabla^2u=0$ をみたす関数 & 8-12 \\
$u_t=\kappa u_{xx}$ & $\dfrac{\partial u}{\partial t}=\alpha\dfrac{\partial^2u}{\partial x^2}$（$\kappa$ はこの本の $\alpha$） & 8-13 \\
陽的差分スキーム、安定条件 $\mu\leqq\frac12$ & 表計算の1歩の式、$r\leqq\frac12$ & 8-13 \\
拡散方程式 & 熱方程式と同じ形の式（もののこさが広がる） & {8-13\\8-14} \\
$\dfrac{\mathrm D\symbf{v}}{\mathrm Dt}$、$\dfrac{d\symbf{v}}{dt}$ & 物質微分 $\dfrac{\mathrm D\vec v}{\mathrm Dt}$ & {8-13\\8-14} \\
\end{longtblr}\par}

%% ---------------- チェック表 ----------------
\bkbacktitle{できるようになったことチェック表}
{\small 各章の「ふりかえり」を集めた表です。章を学び終えたら1回め、次の巻に進む前に2回めのチェックをしよう。できなかった所は、ページにもどって読みなおします。}\par\vspace{2mm}
%% 自動生成（tools/mkchecklist.py）。各章の「ふりかえり」の項目を集めたもの
\begingroup\footnotesize
\def\bkcb{\tikz[baseline=-0.5ex]\draw[\bkTF, line width=0.7pt, rounded corners=0.5pt] (0,-0.3em) rectangle (0.8em,0.5em);}
\begin{longtblr}{colspec={Q[c,m,wd=9mm] X[l,m] Q[c,m,wd=10mm] Q[c,m,wd=9mm] Q[c,m,wd=9mm]}, width=\linewidth, rowhead=1, colsep=3pt,
  row{1}={bg=bkVol!14, font=\sffamily\bfseries\scriptsize}, hlines={0.4pt,bkLine}, rowsep=0.7pt}
章 & できるようになったこと & ページ & 1回め & 2回め \\
{\sffamily\bfseries\color{\bkTT}8-1} & 地図の等高線や天気図の等圧線から、高い所・低い所・急な所を読みとれる。 & \pageref{bk:8-1:kankaku} & \bkcb & \bkcb \\
 & 2つの変数の関数を、表・鳥かん図・等高線の3通りで表し、表計算と Python で作れる。 & \pageref{bk:8-1:hyou} & \bkcb & \bkcb \\
 & 等高線の式と切り口の関数を求め、3つの変数の関数の等値面を説明できる。 & \pageref{bk:8-1:daen} & \bkcb & \bkcb \\
{\sffamily\bfseries\color{\bkTT}8-2} & 等高線や数の表から、東の坂と北の坂を別々に読みとれる。 & \pageref{bk:8-2:hyou} & \bkcb & \bkcb \\
 & ほかを止めて1つだけ動かす偏微分を計算し、点での値の意味を言える。 & \pageref{bk:8-2:henbibun} & \bkcb & \bkcb \\
 & 2回の偏微分を計算して順番によらないことを確かめ、3つの変数の量に使える。 & \pageref{bk:8-2:nikai} & \bkcb & \bkcb \\
{\sffamily\bfseries\color{\bkTT}8-3} & ななめの変化を、東の変化と北の変化のたし算で見積もれる。 & \pageref{bk:8-3:kaidan} & \bkcb & \bkcb \\
 & 全微分で、近似の値と誤差の伝わり方を計算できる。 & \pageref{bk:8-3:zenbibun} & \bkcb & \bkcb \\
 & 連鎖律とヤコビ行列で、手の動きと面積の倍率（ヤコビアン）を求められる。 & \pageref{bk:8-3:jacobi} & \bkcb & \bkcb \\
{\sffamily\bfseries\color{\bkTT}8-4} & 等高線の地図に、ボールが転がりはじめる向き（等高線と直角）をかきこめる。 & \pageref{bk:8-4:chokkaku} & \bkcb & \bkcb \\
 & 東の坂と北の坂を並べて勾配 grad を作り、どの向きの坂も内積で計算できる。 & \pageref{bk:8-4:houkou} & \bkcb & \bkcb \\
 & 内積の $\cos$ の形から、いちばん急な向きと等高線の向きを説明し、$\vec F=-\mathrm{grad}\,U$ などの物理の式を使える。 & \pageref{bk:8-4:cos} & \bkcb & \bkcb \\
{\sffamily\bfseries\color{\bkTT}8-5} & 場所ごとの矢印を計算して、縮める割合をそろえた矢印の場の図にかける。 & \pageref{bk:8-5:chijimi} & \bkcb & \bkcb \\
 & 主役の6つの場・点電荷の場・勾配の場を見分け、流線（力線）の式を求められる。 & \pageref{bk:8-5:ryusenshiki} & \bkcb & \bkcb \\
 & 座標の軸を回して見たときの矢印の成分を、回転行列の転置で計算できる。 & \pageref{bk:8-5:kaiten} & \bkcb & \bkcb \\
{\sffamily\bfseries\color{\bkTT}8-6} & 板を小さなタイルに分けて重さを表計算でたし、重積分の記号を「何を・どこで・1枚の量」で読める。 & \pageref{bk:8-6:tile} & \bkcb & \bkcb \\
 & 長方形の範囲の重積分を、ふつうの定積分を2回くり返す形に直して計算できる。 & \pageref{bk:8-6:naosu} & \bkcb & \bkcb \\
 & 場合分けの表を使って、三角形・円・立体の重積分を、正しいふつうの定積分に直して計算できる。 & \pageref{bk:8-6:basho} & \bkcb & \bkcb \\
{\sffamily\bfseries\color{\bkTT}8-7} & 道を小さな1歩に分け、力と1歩の内積をたして仕事を表計算で求め、線積分の記号を読める。 & \pageref{bk:8-7:ippo} & \bkcb & \bkcb \\
 & 軸に平行な道（折れ線）の線積分を、ふつうの定積分に直して計算できる。 & \pageref{bk:8-7:jiku} & \bkcb & \bkcb \\
 & 場合分けの表を使って、ななめ・曲線・円の道、保存力、閉じた道の線積分を計算できる。 & \pageref{bk:8-7:basho} & \bkcb & \bkcb \\
{\sffamily\bfseries\color{\bkTT}8-8} & 網を通る水の量が、流れの「面に直角な成分」×面積で決まることを説明できる。 & \pageref{bk:8-8:chokkaku} & \bkcb & \bkcb \\
 & 平らな面の流束を、法線ベクトルの内積と重積分で求められる。 & \pageref{bk:8-8:taira} & \bkcb & \bkcb \\
 & 場合分けの表で、曲がった面や閉じた面の面積分をふつうの定積分に直せる。 & \pageref{bk:8-8:basho} & \bkcb & \bkcb \\
{\sffamily\bfseries\color{\bkTT}8-9} & 矢印の図にマスを置いて、わき出しと吸いこみを見分けられる。 & \pageref{bk:8-9:hako} & \bkcb & \bkcb \\
 & マスの辺を通る量から、表計算で発散を求められる。 & \pageref{bk:8-9:masu} & \bkcb & \bkcb \\
 & 発散を偏微分の和で計算し、流束や連続の式と結びつけられる。 & \pageref{bk:8-9:shiki} & \bkcb & \bkcb \\
{\sffamily\bfseries\color{\bkTT}8-10} & 矢印の場の図から、小さな風車が回るかどうかと、その向きを見分けられる。 & \pageref{bk:8-10:fusha} & \bkcb & \bkcb \\
 & マスのまわりを1周してたす量（循環）を表計算で求め、面積でわって回転を見積もれる。 & \pageref{bk:8-10:masu} & \bkcb & \bkcb \\
 & 平面と3次元の rot を偏微分で計算し、保存力の見分けや、電流のまわりの磁場に使える。 & \pageref{bk:8-10:heimen} & \bkcb & \bkcb \\
{\sffamily\bfseries\color{\bkTT}8-11} & マスごとのわき出し（うず）を全部たすと、外のふちから出る量（外の1周）になることを、表計算で確かめられる。 & \pageref{bk:8-11:uchikeshi} & \bkcb & \bkcb \\
 & 閉じた面や閉じた道の積分を、2つの定理で「計算が楽な側」の積分に取りかえて求められる。 & \pageref{bk:8-11:gauss} & \bkcb & \bkcb \\
 & 対称性に合った袋や道を選び、ガウスの法則・アンペールの法則で、場の大きさを求められる。 & \pageref{bk:8-11:hosoku} & \bkcb & \bkcb \\
{\sffamily\bfseries\color{\bkTT}8-12} & マスの値と「まわりの平均」との差を出し、温まる所と冷える所を見分けられる。 & \pageref{bk:8-12:sa} & \bkcb & \bkcb \\
 & 「どのマスもまわりの平均と同じ」のくり返しで、板の落ち着いた温度を求められる。 & \pageref{bk:8-12:ita} & \bkcb & \bkcb \\
 & ラプラシアンを偏微分で計算し、球や円柱の形の場合にも使える。 & \pageref{bk:8-12:lap} & \bkcb & \bkcb \\
{\sffamily\bfseries\color{\bkTT}8-13} & 棒の温度を「となりとの差」の表計算で計算できる。 & \pageref{bk:8-13:masu} & \bkcb & \bkcb \\
 & こわれない条件 $r\leqq\frac12$ を1歩の倍率で説明できる。 & \pageref{bk:8-13:joken} & \bkcb & \bkcb \\
 & 熱と波の方程式を作り、代入して確かめられる。 & \pageref{bk:8-13:netsu} & \bkcb & \bkcb \\
{\sffamily\bfseries\color{\bkTT}8-14} & 4つの式を、「わき出し」と「うず」のことばで読める。 & \pageref{bk:8-14:yomu} & \bkcb & \bkcb \\
 & 積分形と微分形のつながり、光の速さの出し方を言える。 & \pageref{bk:8-14:sekibun} & \bkcb & \bkcb \\
 & 流れの式・熱の式の項を、この巻のことばに訳せる。 & \pageref{bk:8-14:ns} & \bkcb & \bkcb \\
\end{longtblr}
\endgroup


%% ---------------- つぎの巻へ（第9巻の予告と、この巻の道具がこの先どこで使われるか） ----------------
\bkbacktitle{つぎの巻へ}
\newcommand\bknextpill[1]{\bk@pill{#1}}
\def\bknextrow{1}
\newcommand\bknextcard[4]{%
  \begin{tcolorbox}[enhanced, nobeforeafter, box align=top, width=0.485\linewidth, colframe=bkVol9, colback=white, boxrule=0.9pt, arc=2.5mm,
    left=2.5mm, right=2.5mm, top=1.5mm, bottom=2mm, equal height group=bknext\bknextrow]
  \raggedright\bknextpill{#1}\enspace{\sffamily\bfseries\small #2}\par\vspace{1mm}
  {\sffamily\bfseries\footnotesize\color{bkDeep9}#3}\par\vspace{1mm}
  {\footnotesize\setlength{\parindent}{0pt}{\sffamily\bfseries\color{bkInk}使う道具}\enspace #4}
  \end{tcolorbox}}
空間に広がる量を、微分と積分で読みとく道具がそろいました。第9巻「波と変化を読みとく」では、波や音を $\sin$ と $\cos$ に分け、微分方程式をかけ算で解きます。8-13 の熱方程式と波動方程式が、その入口です。\par\vspace{2.5mm}
\noindent\bknextcard{9-0}{音を分ける、変化を予想する}{声を成分に分けると、何が見える？}{波動方程式と、その答え $f(x-ct)$（8-13）。光も電磁波として伝わる（8-14）。}\hfill
\bknextcard{9-1}{波をたし合わせて形を作る：フーリエ級数}{sinをたしていくと、四角い波ができる}{$\sin$ の形の温度は、形を保ったまま小さくなる（8-13）。弦のゆれも $\sin$ の形。}\par\vspace{1.2mm}
\def\bknextrow{2}\noindent\bknextcard{9-2}{成分の取り出し方：直交性}{かけて積分すると、ほしい成分だけが残る}{対称性で0になる積分（8-6 の場合G）と、内積（4-12）。}\hfill
\bknextcard{9-3}{複素数で書くフーリエ級数}{sinとcosを、1つの式でまとめて書く}{熱方程式（8-13）を、変数分離とフーリエ級数で「分けて解いて、たし直す」。}\par\vspace{1.2mm}
\def\bknextrow{3}\noindent\bknextcard{9-4}{スペクトル：フーリエ変換}{くり返さない音も、周波数に分けられる}{$-\infty$ から $\infty$ までの積分と、ガウス積分 $\int e^{-x^2}\,\dd x=\sqrt\pi$（8-6）。}\hfill
\bknextcard{9-11}{行列の指数関数と固有値}{固有値が、落ち着くかどうかを決める}{熱の1歩の行列の固有値と、こわれない条件 $r\leqq\frac12$（8-13）。}\par

\clearpage
\bkbkhead{この巻の道具は、この先でこう使う}
{\small 第8巻の章と、その道具を使う第9巻から先の章です（章のはじめの「道しるべ」の「この章の先に」から集めました）。}\par\vspace{3mm}
%% 自動生成（tools/mknext.py）。各章の道しるべ「この章の先に」から、後ろの巻の章を集めたもの
\begin{longtblr}{colspec={Q[l,m,wd=58mm] X[l,m]}, width=\linewidth, rowhead=1, row{1}={bg=bkVol!14, font=\sffamily\bfseries\footnotesize},
  hlines={0.4pt,bkLine}, rowsep=1.5pt, cells={font=\footnotesize}}
この巻の章 & この道具を使う章 \\
{\bknextpill{8-13}\ 熱と波の伝わり方：偏微分方程式} & {\bknextpill{9-1}\ 波をたし合わせて形を作る：フーリエ級数} \\
{\bknextpill{8-14}\ div と rot で書かれた式を読もう} & {\bknextpill{9-0}\ 音を分ける、変化を予想する} \\
\end{longtblr}

\vspace{4mm}
\bkfukidashi{delta}{grad・div・rot と積分の地図は、物理シリーズの電磁気・流体・熱で、ずっと使うよ。自信のない章を見なおしてから、第9巻へ進もう！}

%% ---------------- さくいん ----------------
\bkbacktitle{さくいん}
{\small 数字はページです。その語がはじめて出てくるページや、くわしく説明しているページをのせています。}\par\vspace{2mm}
\begin{multicols}{2}\small\setlength{\parskip}{0pt}
\directlua{bk.printidx()}
\end{multicols}

%% ---------------- 奥付 ----------------
\bkpadfour
\global\bk@decofalse\bkfolio{false}
\null\vfill
\noindent\begin{minipage}{\linewidth}\small
{\sffamily\bfseries\Large\color{bkDeep}物理を理解するための数学　第8巻}\par\vspace{1mm}
{\sffamily\bfseries\large 空間に広がる量}\par\vspace{3mm}
{\color{bkVol}\hrule height 1pt}\vspace{3mm}
\begin{tabular}{@{}l@{\hspace{6mm}}l@{}}
編者 & 「物理を理解するための数学」編集委員会 \\
 & 委員長　小学校長（まとめ役） \\
 & 委員　数学・電気工学・機械工学・教育学・量子力学・建築工学・\\
 & \phantom{委員　}半導体工学の大学教授、幼稚園教諭、塾講師、\\
 & \phantom{委員　}特別支援学校教諭、心理士、科学解説者、小学校教諭 \\
デザイン & ユニバーサルデザイン担当委員 \\
挿絵 & イラスト担当委員 \\
版 & 第1版（第14回編集委員会）　2026年10月 \\
\end{tabular}\par\vspace{3mm}
{\color{bkVol}\hrule height 1pt}\vspace{2mm}
{\footnotesize この本の問題の答えは、巻末の「こたえ」にあります。くわしい解き方は、別冊の問題集と指導書にのせます。}
\end{minipage}
\makeatother
